\documentclass{exam}

\newif\ifA
\newif\ifkey
\newif\ifprintselected   % required by exam_config's \ansbox

%%%%%%%%%%%%%%%%%%%%%%%%%%%%%%%%%%%%%%%%%%%%%%%%%%%%%%
%%% SET THESE IF VARIABLES
%%%%%%%%%%%%%%%%%%%%%%%%%%%%%%%%%%%%%%%%%%%%%%%%%%%%%%
% Lab name -- labs are named by the lesson they belong to, never by a
% sequential lab number (Neil, 2026-10-12): this is the L16 Lab.
\def\labLesson{16}

% is this the key or not?
\ifdefined\iskey
	\ifnum\iskey=1
		\keytrue
	\else
		\keyfalse
	\fi
\else
	%% Pick one manually
	\keyfalse % if this is NOT the key
%	\keytrue     % if this IS the key
\fi

% set up the header and footer - change for every term
\ifdefined\thisTermIn
	\def\thisterm{\thisTermIn}
\else
	\def\thisterm{Spring 2027}
\fi
%
\def\thisclass{ECE 448}

%%%%%%%%%%%%%%%%%%%%%%%%%%%%%%%%%%%%%%%%%%%%%%%%%%%%%%
%%% END SET THESE IF VARIABLES
%%%%%%%%%%%%%%%%%%%%%%%%%%%%%%%%%%%%%%%%%%%%%%%%%%%%%%
% packages and shortcuts
\input{myPackages_gr}
\usepackage[hidelinks, linkcolor=red]{hyperref}
\input{myShortcuts}
\setlength{\parindent}{0pt}
\setlength{\parskip}{6pt}

%% exam class customization
\input{exam_config}
% ECE 448 document fonts.
%
% Body text is Barlow (SIL OFL; the faces and the license are in
% latex/fonts/Barlow/). Barlow is not in TeX Live, so it is loaded as an
% OpenType face through fontspec -- which means the document must be built
% with lualatex, not pdflatex. build_practice.sh and build_lab.sh both do.
%
% Compiling with pdflatex anyway is not an error: the guard below falls back
% to the LaTeX default so the source still compiles under any engine and can
% be syntax-checked. The resulting PDF just is not the one we ship.
%
% Paths are relative to latex/, which is where both build scripts run
% (each does cd "$(dirname "$0")"). Loading by path rather than by installed
% family name means no system font install and no fontconfig, so a fresh
% clone builds a byte-comparable PDF anywhere.
%
% MATH IS DELIBERATELY LEFT IN COMPUTER MODERN. Barlow has no math companion,
% and CM math under sans body text is already the house pairing: the decks set
% words in Source Sans Pro and let MathJax render symbols in its CM-derived
% face. Moving math to a sans as well would mean unicode-math plus a sans math
% font (Fira Math being the plausible one) and re-checking every equation in
% the course for symbol coverage. That is a separate decision, not a side
% effect of this one.

\usepackage{iftex}
\ifLuaTeX
	\usepackage{fontspec}
	\defaultfontfeatures{Ligatures=TeX}
	\setmainfont{Barlow}[
		Path           = fonts/Barlow/,
		Extension      = .ttf,
		UprightFont    = *-Regular,
		ItalicFont     = *-Italic,
		BoldFont       = *-Bold,
		BoldItalicFont = *-BoldItalic,
	]
	\setsansfont{Barlow}[
		Path           = fonts/Barlow/,
		Extension      = .ttf,
		UprightFont    = *-Regular,
		ItalicFont     = *-Italic,
		BoldFont       = *-SemiBold,
		BoldItalicFont = *-BoldItalic,
	]
	% microtype earns its place here specifically because of Barlow. Its
	% glyphs are narrower than Computer Modern's, so justified prose sets
	% tighter and TeX has to squeeze interword space to fit; two paragraphs
	% in the L02/L03 sets went overfull on the first Barlow build. Protrusion
	% and font expansion absorb that -- both are lualatex features, so this
	% belongs inside the guard with everything else.
	\usepackage{microtype}

	% Mono stays Latin Modern: Barlow ships no monospaced face, and verbatim
	% and code need one that actually is fixed-pitch.
\fi
   % Barlow body text (needs lualatex)
\printselectedfalse      % full worked solutions in the key, not answers-only
\CorrectChoiceEmphasis{\color{red}}
\SolutionEmphasis{\color{red}}
\renewcommand\questionlabel{\textbf{\thequestion.}}
\renewcommand{\thepartno}{\alph{partno}}
\renewcommand{\partlabel}{(\thepartno)}

% Fillable table cell: blank in the student copy, the expected value in red on
% the KEY. Fixed-width centered columns keep the blank cells writable.
\newcolumntype{K}[1]{>{\centering\arraybackslash}p{#1}}
\newcommand{\kv}[1]{\ifkey{\color{red}#1}\fi\rule[-0.10in]{0pt}{0.34in}}

% Answers are not in this public file. Every solution, \kv value, and \ansbox
% value is \keysol{<id>}, defined in the private ECE448-faculty repo at
% lab-keys/src/ECE448_Lab_L16_FHSS_solutions.tex; build_lab.sh puts that
% folder on TEXINPUTS (KEYDIR) for the KEY. A KEY build without it stops.
\long\def\defkeysol#1#2{\expandafter\long\expandafter\def\csname keysol@#1\endcsname{#2}}
\def\keysol#1{\ifcsname keysol@#1\endcsname\csname keysol@#1\endcsname\fi}
\ifkey
	\IfFileExists{ECE448_Lab_L16_FHSS_solutions.tex}%
		{\input{ECE448_Lab_L16_FHSS_solutions}}%
		{\errmessage{KEY build needs ECE448_Lab_L16_FHSS_solutions.tex from ECE448-faculty/lab-keys/src (set KEYDIR)}}
\fi

% print the answers if this is the key
\ifkey
	\printanswers % if this IS the key
\else
	\noprintanswers  %if this is NOT the key
\fi

% print the header and footer
% need to print key on top if this is the key
\ifkey
	\firstpageheader{}{\color{red}{ KEY}}{}
	\runningheader{\thisclass\hspace{1pt} -- \thisterm }{L\labLesson\ Lab\hspace{1pt}\color{red}{ KEY}}{\thepage\hspace{1pt} of\hspace{1pt} \numpages}
	\firstpagefooter{}{}{}
	\runningfooter{}{}{}
\else
	\firstpageheader{}{}{}
	\runningheader{\thisclass \hspace{1pt} -- \thisterm }{L\labLesson\ Lab\hspace{1pt}}{\thepage\hspace{1pt} of\hspace{1pt} \numpages}
	\firstpagefooter{}{}{}
	\runningfooter{}{}{}
\fi

\runningheadrule


\begin{document}
\pagestyle{headandfoot}   % exam class

%%%%%%%%%%%%%%%%%%%%%%%%%%%%%%%%%%%%%%%%%%%%%%%%%%%%%%
% title block
%%%%%%%%%%%%%%%%%%%%%%%%%%%%%%%%%%%%%%%%%%%%%%%%%%%%%%

\begin{center}
USAF ACADEMY \\
DEPARTMENT OF ELECTRICAL AND COMPUTER ENGINEERING \\
LESSON \labLesson\ LAB - FREQUENCY HOPPING \\
\textbf{Lab Packet} \\
\thisterm \\[4mm]
\end{center}

\vspace{-4pt}
\textbf{Name:} \rule[-2pt]{2.6in}{0.4pt} \hfill
\textbf{Section:} \rule[-2pt]{0.9in}{0.4pt} \hfill
\textbf{Date:} \rule[-2pt]{1.1in}{0.4pt}

\vspace{4pt}

\fullwidth{\textbf{Learning Objective 6.5: } I can implement FHSS transmission and
reception, including hop-sequence synchronization.}

\vspace{2pt}

This packet is the L\labLesson\ Lab, in three parts, and you turn it in once, complete,
before Lesson 17. Do \textbf{Part 1} on your own before the lab period: run the link
flowgraph and work out the settings the bench uses. \textbf{Part 2} is what you
build and observe, recorded in class: Questions 3 to 5 on your own laptop, and
Question 6 at a PlutoSDR station with your team. Do \textbf{Part 3} on your own after
the lab: explain, in words, what changed and why.

\textbf{Grading.} You submit the packet once, we grade it against Objective 6.5 on
the course's 1 / 0 scale, and you may revise it under course policy. Paste screenshots where a
question asks for one.

\textbf{Transmitting.} In class we transmit only into a cable, into a shielded
enclosure, or at very low power in an ISM band, and only when the lesson says to.
Questions 1 to 5 run on your laptop and transmit nothing. Question 6 transmits from a
PlutoSDR's TX port, through a cable with an inline 30 dB SMA attenuator, to the same
Pluto's RX port, at 915 MHz. Connect the cable and attenuator before you run the
flowgraph, and never connect an antenna in this lab. Never set TX attenuation below
30 dB or change the slider's range; the RX input is rated to only a few dBm. We have
run the PlutoSDR flowgraph only in simulation, so Question 6 is its first run on
hardware.

The flowgraphs are on the L16 lesson page, \texttt{L16-fhss-tx.grc} and
\texttt{L16-fhss-link.grc}, and on the L16 Lab page, \texttt{L16-fhss-pluto.grc}. Every
number below assumes their settings: Lesson 8's FSK at 1000 baud, sampled at 128 kHz;
one packet per hop, made of \texttt{guard} guard bytes, the 12-byte access code and
header, the 14-byte payload \texttt{USAFA ECE 448} with its newline, and
\texttt{guard} more guard bytes; $N$ channels 10 kHz apart; and key 448, whose sync
channel for $N = 8$ is channel 7.

\begin{questions}

\fullwidth{\large\textbf{Part 1: Pre-Lab} \normalsize(individual, before the lab period)}

%%%%%%%%%%%%%%%%%%%%%%%%%%%%%%%%%%%%%%%%%%%%%%%%%%%
\question \textbf{Run the link.} Open \texttt{L16-fhss-link.grc} in GRC and run it
with every slider at its default.

\begin{parts}

	\begin{minipage}[t]{0.58\linewidth}
		\part Once the Receiver label reads ``locked,'' wait 30 s and record the Receiver
		and Packets labels.
		\begin{solution}[0.7in]\keysol{sol1}\end{solution}
	\end{minipage}\hfill%
	\begin{minipage}[t]{0.37\linewidth}
		\raggedleft \vspace{12pt}
		\ansbox[1.4in]{\keysol{ans1}}
	\end{minipage}

	\part Paste a screenshot of the running flowgraph's window.
		\begin{solution}[1.6in]\keysol{sol2}\end{solution}

\end{parts}

%%%%%%%%%%%%%%%%%%%%%%%%%%%%%%%%%%%%%%%%%%%%%%%%%%%
\newpage
\question \textbf{Settings.} Work out the numbers the flowgraphs use.

\begin{parts}

	\part Set \texttt{guard} to 4 bytes. Find the bits in one hop, the time on one hop,
	the hop rate, the time for one cycle of 8 channels, and the samples in one hop at
	128 kHz.

	\begin{center}
	\renewcommand{\arraystretch}{1.25}
	\begin{tabular}{| l | K{0.75in} | K{0.75in} | K{0.75in} | K{0.75in} | K{0.85in} |}
	\hline
	& Bits per hop & Hop time & Hop rate & Cycle & Samples \\\hline
	\texttt{guard} = 4 & \kv{\keysol{kv1}} & \kv{\keysol{kv2}} & \kv{\keysol{kv3}} & \kv{\keysol{kv4}} & \kv{\keysol{kv5}} \\\hline
	\end{tabular}
	\end{center}

		\begin{solution}[0.9in]\keysol{sol3}\end{solution}

	\begin{minipage}[t]{0.58\linewidth}
		\part With \texttt{guard} = 4, how long does each guard last?
		\begin{solution}[0.6in]\keysol{sol4}\end{solution}
	\end{minipage}\hfill%
	\begin{minipage}[t]{0.37\linewidth}
		\raggedleft \vspace{12pt}
		\ansbox[1.1in]{\keysol{ans2}}
	\end{minipage}

	\part For the hopper you build in Question 3, choose your own order of the 8
	channels, using each once, and write the hop frequencies in that order. With
	\texttt{dwell} = 32{,}000 samples at 128 kHz, find the time on one hop and the time
	for one cycle.
		\begin{solution}[1.0in]\keysol{sol5}\end{solution}

\end{parts}

%%%%%%%%%%%%%%%%%%%%%%%%%%%%%%%%%%%%%%%%%%%%%%%%%%%
\newpage
\fullwidth{\large\textbf{Part 2: Measurements} \normalsize(in class)}

\question \textbf{Build a hopper and a de-hopper.} Start GRC with an empty canvas.
Do not open \texttt{L16-fhss-tx.grc}: build the hopper and the de-hopper yourself as
the lab page lists, with your order from Question 2(c).

\begin{parts}

	\part Paste a screenshot of your running flowgraph's window with no Delay block.
	Mark where the jammer appears on each waterfall.
		\begin{solution}[1.6in]\keysol{sol6}\end{solution}

	\part Paste a screenshot with the Delay block set to \texttt{dwell}. Explain what
	the De-hopped waterfall shows.
		\begin{solution}[1.4in]\keysol{sol7}\end{solution}

	\part Describe the De-hopped waterfall with the delay at 2000 samples.
		\begin{solution}[1.0in]\keysol{sol8}\end{solution}

\end{parts}

%%%%%%%%%%%%%%%%%%%%%%%%%%%%%%%%%%%%%%%%%%%%%%%%%%%
\newpage
\question \textbf{Jamming.} Follow the lab page's steps for \texttt{L16-fhss-link.grc}.
Record the Packets label after 30 s at each setting, as correct of sent, and as a
percent.

\begin{parts}

	\part Fill in the table.

	\begin{center}
	\renewcommand{\arraystretch}{1.25}
	\begin{tabular}{| l | K{0.65in} | K{0.65in} | K{0.65in} | K{0.65in} | K{0.65in} |}
	\hline
	\texttt{nch} & 8 & 1 & 8 & 2 & 4 \\\hline
	Jammer & 2, ch 4 & 0.7 & 0.7, ch 4 & 2, ch 1 & 2, ch 1 \\\hline
	Delivered & \kv{\keysol{kv6}} & \kv{\keysol{kv7}} & \kv{\keysol{kv8}} & \kv{\keysol{kv9}} & \kv{\keysol{kv10}} \\\hline
	\end{tabular}
	\end{center}

		\begin{solution}[0.9in]\keysol{sol9}\end{solution}

	\begin{minipage}[t]{0.58\linewidth}
		\part Compare the two results at jammer amplitude 0.7, and explain in words why the
		hopping link lost so many fewer packets.
		\begin{solution}[1.0in]\keysol{sol10}\end{solution}
	\end{minipage}\hfill%
	\begin{minipage}[t]{0.37\linewidth}
		\raggedleft \vspace{12pt}
		\ansbox[1.3in]{\keysol{ans3}}
	\end{minipage}

	\part Explain in words why 2 channels lost about half the packets and 4 channels about
	a quarter.
		\begin{solution}[0.8in]\keysol{sol11}\end{solution}

	\part Record the Receiver and Packets labels with the jammer at 3 on channel 7 turned
	on after the receiver locked, and again set as the slider defaults before the run.
		\begin{solution}[0.9in]\keysol{sol12}\end{solution}

\end{parts}

%%%%%%%%%%%%%%%%%%%%%%%%%%%%%%%%%%%%%%%%%%%%%%%%%%%
\newpage
\question \textbf{Sync.} Follow the lab page's steps.

\begin{parts}

	\begin{minipage}[t]{0.58\linewidth}
		\part Record the Packets label with the hop offset at $+1$ and at $-1$.
		\begin{solution}[0.6in]\keysol{sol13}\end{solution}
	\end{minipage}\hfill%
	\begin{minipage}[t]{0.37\linewidth}
		\raggedleft \vspace{12pt}
		\ansbox[1.1in]{\keysol{ans4}}
	\end{minipage}

	\part With \texttt{guard} = 4, fill in the table, and compare the edges with the guard time from Question 2(b).

	\begin{center}
	\renewcommand{\arraystretch}{1.25}
	\begin{tabular}{| l | K{0.55in} | K{0.55in} | K{0.55in} | K{0.55in} | K{0.55in} | K{0.65in} |}
	\hline
	Timing (ms) & $-24$ & $+24$ & $-30$ & $+30$ & $-36$ & $+36$ \\\hline
	Delivered & \kv{\keysol{kv11}} & \kv{\keysol{kv12}} & \kv{\keysol{kv13}} & \kv{\keysol{kv14}} & \kv{\keysol{kv15}} & \kv{\keysol{kv16}} \\\hline
	\end{tabular}
	\end{center}

		\begin{solution}[1.0in]\keysol{sol14}\end{solution}

	\begin{minipage}[t]{0.58\linewidth}
		\part Record the three times from moving the path delay to the first correct packet.
		\begin{solution}[0.9in]\keysol{sol15}\end{solution}
	\end{minipage}\hfill%
	\begin{minipage}[t]{0.37\linewidth}
		\raggedleft \vspace{12pt}
		\ansbox[1.1in]{\keysol{ans5}}
	\end{minipage}

\end{parts}

%%%%%%%%%%%%%%%%%%%%%%%%%%%%%%%%%%%%%%%%%%%%%%%%%%%
\newpage
\question \textbf{Frequency hopping on a PlutoSDR.} At a PlutoSDR station, with your
team, follow the lab page's steps: connect the Pluto's TX port to its RX port through
the 30 dB attenuator, check the Pluto with \texttt{iio\_info -s}, and run
\texttt{L16-fhss-pluto.grc}. When you finish, stop the flowgraph and unplug the Pluto
before you disconnect the cable.

\begin{parts}

	\part Once \texttt{USAFA ECE 448} decodes, paste a screenshot that shows both
	waterfalls and the Receiver, Packets, and Received labels.
		\begin{solution}[1.5in]\keysol{sol16}\end{solution}

	\begin{minipage}[t]{0.58\linewidth}
		\part Record the Packets label with the jammer at 2 on channel 4.
		\begin{solution}[0.6in]\keysol{sol17}\end{solution}
	\end{minipage}\hfill%
	\begin{minipage}[t]{0.37\linewidth}
		\raggedleft \vspace{12pt}
		\ansbox[1.1in]{\keysol{ans6}}
	\end{minipage}

	\part Record the Receiver and Packets labels at each TX attenuation, and explain what
	limits the link as the attenuation rises.

	\begin{center}
	\renewcommand{\arraystretch}{1.25}
	\begin{tabular}{| l | K{0.9in} | K{0.9in} | K{0.9in} | K{0.9in} |}
	\hline
	TX attenuation (dB) & 60 & 70 & 80 & 89 \\\hline
	Receiver & & & & \\\hline
	Delivered & & & & \\\hline
	\end{tabular}
	\end{center}

		\begin{solution}[1.0in]\keysol{sol18}\end{solution}

\end{parts}

%%%%%%%%%%%%%%%%%%%%%%%%%%%%%%%%%%%%%%%%%%%%%%%%%%%
\newpage
\fullwidth{\large\textbf{Part 3: Writeup} \normalsize(individual, after the lab)}

\question \textbf{Analysis.}

\begin{parts}

	\part Describe how the packets delivered changed as you changed the number of channels
	and the jammer's amplitude. Use your De-hopped waterfall from Question 3(a) to explain
	why a jammer stronger than the signal takes only the hops on its own channel, and why
	raising its power further does not take more.
		\begin{solution}[1.5in]\keysol{sol19}\end{solution}

	\part Compare the timing errors the link tolerated with 2 guard bytes, from the
	lesson's Timing Error frame, and with 4, from Question 5(b). Explain in words what the
	extra guard bytes take from each hop, and when they are worth it.
		\begin{solution}[1.5in]\keysol{sol20}\end{solution}

	\part From Question 4(d), explain why one jammed channel can keep a receiver from
	ever locking although it takes only 1 hop in 8 from a locked receiver, and propose a
	change to Hop Sync or to the packet that removes that weakness.
		\begin{solution}[1.5in]\keysol{sol21}\end{solution}

\end{parts}

\end{questions}

\vspace{1.0em}
\noindent\textbf{Documentation:} \rule[-2pt]{0.6\linewidth}{0.4pt}

\end{document}
